\section{Metodología}
\label{sec:metodologia}
% Herramientas y versiones, diseño de cada experimento, criterios de análisis, semilla,
% uso de Claude, y el link al repositorio (SOLO acá).

\subsection{Entorno y herramientas}
\label{subsec:entorno}

Todo el trabajo numérico se hizo en Python 3.12.3 sobre Windows 11, con NumPy 1.26.4 (Harris
et al., 2020) para los arreglos y los generadores de números aleatorios, y Matplotlib 3.11.1
(Hunter, 2007) para las figuras. El tipo \texttt{float} de Python corresponde al formato
\texttt{binary64} descrito en la sección \ref{subsec:punto-flotante}, y el tipo \texttt{int} es el
entero de precisión arbitraria de la sección \ref{subsec:enteros}. Las dos únicas dependencias
externas están declaradas en \texttt{requirements.txt} con sus versiones mínimas. El informe se
escribió en \LaTeX{} y se compiló con Tectonic.

Ninguna sumatoria del trabajo usa la función \texttt{sum()} de Python ni las reducciones de NumPy.
Los términos se acumulan siempre con un ciclo \texttt{for} explícito.
El objeto de estudio es el orden en que se encadenan las sumas parciales, y una reducción de
biblioteca puede agrupar los sumandos de una manera que el programa no controla.

\subsection{Organización del código y reproducción de los resultados}
\label{subsec:codigo}

El código está en el directorio \texttt{src/}, con un módulo por experimento
(\texttt{experimento1\_\allowbreak asociativa.py},
\texttt{experimento2\_\allowbreak conmutativa.py} y
\texttt{experimento3\_\allowbreak representaciones.py}), un módulo de utilidades compartidas
(\texttt{comun.py}) y un punto de entrada (\texttt{main.py}). El comando \texttt{python src/main.py}
corre los tres experimentos en unos tres minutos y regenera desde cero las siete figuras y las doce
tablas del informe; pasarle un número corre un solo experimento.

Cada figura se guarda como PDF vectorial en \texttt{informe/figuras/} y cada tabla como cuerpo de
un entorno \texttt{tabular} en \texttt{informe/tablas/}, que el documento incluye con
\texttt{\textbackslash input}. El informe no contiene ningún número que el código no haya escrito
en alguno de esos archivos. La consola imprime solo un resumen de tres líneas por experimento, que
sirve para verificar que la corrida terminó bien; nada de ahí se copia al texto. Todas las figuras
comparten un mismo estilo, fijado una sola vez en \texttt{comun.aplicar\_estilo()}, para que los
tamaños de letra sean legibles en la página impresa.

El repositorio con el código, las figuras, las tablas y las fuentes \LaTeX{} de este informe está
disponible en \url{https://github.com/manuelchouhy/informe-sumatorias}.

\subsection{Diseño de los experimentos}
\label{subsec:diseno}

El primer experimento evalúa $a_N = \sum_{k=1}^{N} (1 + b^{k} - b^{k})$ para
$N = 1, 2, \ldots, 1000$ con las tres asociaciones que fija la consigna,
$1 + b^{k} - b^{k}$, $(1 + b^{k}) - b^{k}$ y $(1 - b^{k}) + b^{k}$, sobre las bases
$b \in \{2, 3, 5, 10\}$ de tipo \texttt{int} y $b \in \{2.0, 3.0, 5.0, 10.0\}$ de tipo
\texttt{float}. Las potencias se calculan multiplicando de a uno y no con el operador
\texttt{**}: con flotantes, \texttt{**} interrumpe el programa con \texttt{OverflowError} al pasar
el mayor valor representable. La multiplicación sucesiva devuelve $\infty$, que es el comportamiento
que el experimento necesita observar. Los desbordes no se atrapan; se dejan llegar a
$\infty$ y a \texttt{NaN}, se grafican tal como salen y se explican en la discusión. Las preguntas
adicionales de la consigna se resuelven con el mismo código y otros parámetros: bases
$b \in \{0.0,\, 0.1,\, 0.5,\, 0.9,\, 1.0\}$ para el régimen $b \le 1$, bases
$b \in \{-2.0, -3.0, -5.0, -10.0\}$ para el signo negativo, una simulación de la aritmética de
enteros con signo de 64 bits para el caso de ancho fijo, y una medición de tiempo sobre
$N \in \{1000,\, 5000,\, 10\,000,\, 20\,000,\, 50\,000,\, 100\,000\}$ para el costo de los enteros
de precisión arbitraria. La simulación de 64 bits reduce cada operación intermedia módulo $2^{64}$
en complemento a dos, que es la convención de Java y la del tipo \texttt{int64} de NumPy. Se la
implementó de forma explícita en lugar de usar \texttt{int64} porque cada multiplicación que
desborda emite un aviso de NumPy en la consola, y el resultado es el mismo.

El segundo experimento evalúa $b_N = \sum_{k=1}^{N} 1/(k(k+1))$ con cuatro algoritmos de suma:
recorriendo $k$ de 1 a $N$, que por ser los términos decrecientes equivale a sumar de mayor a menor
módulo; recorriéndolo de $N$ a 1, que es de menor a mayor; sobre una permutación aleatoria de los
índices; y con la suma compensada de Kahan de la ecuación \eqref{eq:kahan}. Los cuatro calculan el
término con la misma expresión: la única diferencia entre ellos es el orden de acumulación, no el
redondeo del término. Las cuatro funciones reciben $N$ como único argumento obligatorio. La de la
permutación acepta además un índice de repetición opcional que selecciona cuál de las permutaciones
derivadas de la semilla maestra se usa, y que solo interviene en el estudio de dispersión. El error relativo se mide contra la forma cerrada
$N/(N+1)$ sobre dos grillas, $N = 10, 20, \ldots, 10\,000$ y
$N = 1000, 2000, \ldots, 1\,000\,000$. La dispersión de la suma aleatoria se estudia repitiéndola
50 veces con $N = 10^{6}$. Los dos órdenes que no dependen de $N$ se calculan con un único barrido
que anota las sumas parciales en cada punto de la grilla; los otros dos se recalculan para cada
$N$, y eso explica casi todo el tiempo de corrida.

Sobre la grilla $N = 1, 10, 20, \ldots, 10\,000$, el tercer experimento compara representaciones
equivalentes. Para $b_N$ enfrenta la suma de $1/(k(k+1))$ contra la telescópica
$\sum (1/k - 1/(k+1))$, las dos medidas contra la forma cerrada. La sucesión
$c_N = \sum_{k=1}^{N} 1/(\sqrt{k^{2}+1} + k)$ no tiene forma cerrada conocida, así que su
representación como cociente se compara contra la representación como diferencia,
$\sqrt{k^{2}+1} - k$, mediante la discrepancia absoluta $D_N$ definida en la sección
\ref{subsec:error-relativo}. Las dos
preguntas adicionales se estudian aparte: las formas cerradas $1 - 1/(N+1)$ y $N/(N+1)$ se comparan
contra el valor racional exacto, calculado con el tipo \texttt{Fraction} de la biblioteca estándar,
porque lo que se mide ahí son los redondeos de las dos expresiones y un flotante de referencia no
serviría de patrón; y los términos individuales se recorren sobre una grilla logarítmica de 500
valores de $k$ hasta $k = 10^{8}$.

\subsection{Criterios de análisis}
\label{subsec:criterios}

La medida de precisión es el error relativo de la ecuación \eqref{eq:error-relativo} siempre que
haya una forma cerrada contra la cual comparar, y la discrepancia absoluta $D_N$ entre
representaciones cuando no la hay. Como el valor de referencia también se evalúa en punto flotante,
el piso de lo medible es del orden de $u$, y un error relativo nulo se lee como coincidencia bit a
bit con esa referencia.

Las figuras de error usan escala logarítmica en el eje vertical porque los valores comparados
abarcan varios órdenes de magnitud. Cuando la serie contiene errores exactamente nulos, que una
escala logarítmica pura no puede representar, se usa escala \texttt{symlog}. Cuando dos curvas se
superponen casi por completo, se dibujan con marcadores en lugar de líneas continuas para que la
superposición se vea en vez de esconderse detrás de la curva de arriba. Junto a los datos se grafican
las cotas teóricas de la sección \ref{sec:marco}, $\sqrt{N}\,u$ de la ecuación \eqref{eq:cota-raiz}
para las sumas y $2k^{2}u$ de la ecuación \eqref{eq:cota-termino} para los términos de $c_N$, de modo
que la comparación entre lo predicho y lo medido quede en la misma figura.

Los tiempos de ejecución que aparecen en el informe son medidas y no valores exactos: cada una es el
mínimo de tres repeticiones tomadas con \texttt{time.perf\_counter()}, que deja afuera las
repeticiones más afectadas por el ruido del sistema operativo. Aun así los valores cambian entre
corridas, y la nota al pie de la tabla correspondiente lo aclara.

Las curvas de los experimentos segundo y tercero dependen de los barridos incrementales, y su
correctitud se verifica dentro del propio código. La función \texttt{\_verificar\_trazas()} de cada
módulo compara bit a bit el resultado del barrido contra la llamada directa a cada función pública,
sobre varios valores de $N$, y detiene la corrida si alguno difiere.

\subsection{Aleatoriedad y reproducibilidad}
\label{subsec:semillas}

El único componente aleatorio del trabajo es la permutación de los índices en la suma randomizada.
Todo lo aleatorio deriva de una semilla maestra fija, declarada en \texttt{comun.py}. La función
\texttt{comun.generador(*claves)} construye un \texttt{numpy.random.default\_rng} a partir de la
secuencia formada por esa semilla seguida de las claves que identifican la corrida. Cada repetición
recibe claves distintas, con lo cual usa un generador distinto e independiente de los demás, y la
corrida completa sigue siendo reproducible.
Reutilizar la misma semilla en cada repetición habría producido 50 permutaciones idénticas y habría
vuelto vacío el estudio de la dispersión.

\subsection{Uso de asistentes de inteligencia artificial}
\label{subsec:ia}

Durante el desarrollo se usó el asistente Claude Code (Anthropic, 2026) como apoyo para la
implementación en Python, la redacción del informe y la revisión de la consistencia entre el texto y
los datos generados. Los autores definieron el diseño de los experimentos, verificaron los
resultados numéricos y son responsables del contenido final. Ningún dato del informe proviene del
asistente: todos salen de las figuras y tablas que produce el código del repositorio.
